\section{习题 3.3：不变量判别与参数分类}
\chr{7}{153}{159}{305}
\K{二次项矩阵 $A$ 的两个特征值决定“同号／异号／有零”；齐次增广矩阵 $H$ 的行列式 $I_3$ 决定是否进一步退化。$I_2=0$ 时要额外检查 $I_3$、非零特征值和线性项。}
\readtip{\textbf{本节计算模板：}先从 $F=ax^2+2hxy+by^2+2gx+2fy+c$ 写出 $I_1=a+b,\ I_2=ab-h^2,\ I_3=\det\begin{psmallmatrix}a&h&g\\h&b&f\\g&f&c\end{psmallmatrix}$。不要把 $xy$ 的\textbf{全部系数}误放到矩阵非对角元，应取其一半。}

\Q{1}{用 $I_1,I_2,I_3$ 判断七条具体二次曲线}
\hint{对原题七个方程各写对称矩阵，算出不变量；无中心且 $I_3=0$ 的项最后用配方检查是否退化。}
\begin{center}\small
\begin{tabular}{@{}c r r r l@{}}\toprule
小问 & $I_1$ & $I_2$ & $I_3$ & 判别\\
\midrule
(1) & $8$ & $-9$ & $63$ & 双曲线 \\
(2) & $8$ & $0$ & $-256$ & 抛物线 \\
(3) & $7$ & $3$ & $-48$ & 实椭圆 \\
(4) & $5$ & $0$ & $0$ & 一条直线（重线） \\
(5) & $-1$ & $-13$ & $-144$ & 双曲线 \\
(6) & $11$ & $15$ & $15$ & 空集 \\
(7) & $5$ & $0$ & $0$ & 两条平行直线 \\
\bottomrule
\end{tabular}
\end{center}
\subpart{(1)—(3) 样例核算} (1) $F=8x^2+6xy-26x-12y+13$，
\[
A=\begin{pmatrix}8&3\\3&0\end{pmatrix},\quad
H=\begin{pmatrix}8&3&-13\\3&0&-6\\-13&-6&13\end{pmatrix}.
\]
由 $\det A=-9<0$、$\det H=63\ne0$ 判双曲线。(2) 的矩阵 $A=\begin{psmallmatrix}4&4\\4&4\end{psmallmatrix}$，$I_2=0$ 且 $I_3=-256\ne0$，故抛物线。(3) $A=\begin{psmallmatrix}3&3\\3&4\end{psmallmatrix}$ 正定，$I_3=-48<0$，故实椭圆。
\subpart{(4)、(7) 必须配方区分} (4) $4x^2+4xy+y^2-4x-2y+1=(2x+y-1)^2$，所以是\textbf{一条直线} $2x+y-1=0$。\par
(7) $x^2-4xy+4y^2-4x+8y+1=(x-2y-2)^2-3$，因此是\textbf{两条平行直线} $x-2y-2=\pm\sqrt3$。
\subpart{(5)、(6) 判断空实} (5) $I_2=-13$，不定的二次型有实双曲线；(6) $I_2=15>0,I_1=11>0$，平移到中心后 $F=Q(X)+I_3/I_2=Q(X)+1>0$，所以无实点。
\ans{七项类型依次为\textbf{双曲线、抛物线、椭圆、一条直线、双曲线、空集、两平行直线}，均与原书答案一致。}

\Q{2}{含参数的二次式什么时候化为两条直线}
\textbf{题意：}$ax^2+4xy+y^2-4x-2y+c=0$。\hint{先完成平方，退化情况一眼能看出来。}
\[
\boxed{F=(2x+y-1)^2+(a-4)x^2+(c-1).}
\]
若 $c=1,a<4$，则 $F=(2x+y-1-\sqrt{4-a}\,x)(2x+y-1+\sqrt{4-a}\,x)$，为两条\textbf{相交}实直线；若 $a=4,c<1$，则 $2x+y-1=\pm\sqrt{1-c}$，为两条\textbf{平行}实直线。
\caution{若把“两条直线”理解为允许重合的两条，$a=4,c=1$ 还给出 $F=(2x+y-1)^2=0$，是一条重线。原书仅列“两条不同直线”的两种情况；本讲义明确重线这个边界情形。}
\ans{两条相交线条件 $a<4,c=1$；两条平行线条件 $a=4,c<1$。}

\Q{3}{两个含参数 $t$ 的方程如何逐区间分类}
\subpart{(1) $F=(1+t^2)(x^2+y^2)-4txy+2t(x+y)+2=0$}\par
\hint{旋转 $45^\circ$，令 $u=(x+y)/\sqrt2,v=(x-y)/\sqrt2$；分离平方项。}
\[
F=(t-1)^2u^2+(t+1)^2v^2+2\sqrt2\,tu+2.
\]
当 $t\ne\pm1$ 配方：
\[
(t-1)^2\left(u+\frac{\sqrt2t}{(t-1)^2}\right)^2+(t+1)^2v^2
=\frac{4t-2}{(t-1)^2}.
\]
故 $t<1/2$ 时空集；$t=1/2$ 时是一个点；$t>1/2$ 且 $t\ne1$ 时为椭圆。特例 $t=1$ 代回得 $4v^2+2\sqrt2u+2=0$，为真抛物线。$t=-1$ 时 $4u^2-2\sqrt2u+2=0$ 判别式小于零，仍为空集。
\ans{$t<1/2$：空集；$t=1/2$：一点；$t=1$：抛物线；$t>1/2$ 且 $t\ne1$：椭圆。}
\subpart{(2) $F=x^2-4xy+4y^2-4txy-2tx+8y+3-2t=0$}\par
\hint{计算完整不变量并在 $I_2,I_3$ 为零的参数上逐个单独验证。}
\[
I_1=5,\quad I_2=-4t(t+2),\quad I_3=8(t+2)(t+1)(t-1).
\]
据符号和边界值，可得以下\textbf{无遗漏分类}：
\begin{center}\small
\begin{tabular}{@{}ll@{}}\toprule
参数 & 曲线类型\\
\midrule
$t<-2$ & 双曲线 \\
$t=-2$ & 空集\\
$-2<t<-1$ & 空集\\
$t=-1$ & 一个点\\
$-1<t<0$ & 椭圆\\
$t=0$ & 抛物线\\
$0<t<1$ & 双曲线\\
$t=1$ & 两条相交直线\\
$t>1$ & 双曲线\\\bottomrule
\end{tabular}
\end{center}
例如 $t=1$ 时令 $X=x-1,Y=y$，式子变为 $X^2-8XY+4Y^2=0$，即两条直线 $X=(4\pm2\sqrt3)Y$；当 $t=-1$，$A$ 正定而中心处 $F=0$，所以只剩一个点。\ans{两小题参数区间与原书一致；每个临界值已单独处理。}

\Q{4}{证明 $I_2=0$ 时 $I_1I_3\le0$}
\hint{秩 1 的对称矩阵正交对角化后变为 $\operatorname{diag}(\lambda,0)$，其中 $\lambda=I_1$。}
正交坐标下 $F=\lambda u^2+2p u+2q v+r$，增广矩阵
\[
\begin{pmatrix}\lambda&0&p\\0&0&q\\p&q&r\end{pmatrix}
\]
的行列式为 $I_3=-\lambda q^2$；于是 $I_1I_3=-\lambda^2q^2\le0$。若 $A=0$，则 $I_1=I_2=0$，亦自动成立。
\ans{得证；正交变换不改变 $I_1,I_3$。}

\Q{5}{抛物线为什么必须满足 $I_1I_3<0$}
\hint{真抛物线的正交标准形为 $\lambda u^2+2qv=0$，其中 $\lambda q\ne0$。}
标准形的 $I_1=\lambda$、$I_3=-\lambda q^2$，所以 $I_1I_3=-\lambda^2q^2<0$。反过来，若 $I_2=0$ 且 $I_1I_3<0$，则 $\lambda\ne0,q\ne0$，配方和平移给真正的抛物线。
\ans{题目要求的必要性得证；再加 $I_2=0$ 还能构成充要判别条件。}

\Q{6}{圆的充要条件}
\textbf{题意：}证明在直角坐标系中 $\Gamma:F=0$ 是实圆当且仅当 $I_1^2=4I_2$ 且 $I_1I_3<0$。\hint{第一条件迫使二次项矩阵为标量矩阵。}
\[
I_1^2-4I_2=(a_{11}-a_{22})^2+4a_{12}^2=0
\]
得 $a_{11}=a_{22}=\lambda$、$a_{12}=0$。第二条件要求 $\lambda\ne0$，且配方后的半径平方
\[
r^2=\frac{b_1^2+b_2^2-\lambda c}{\lambda^2}
=-\frac{I_3}{\lambda I_2}>0.
\]
所以标准形 $(x+h)^2+(y+k)^2=r^2$ 有实圆。反向计算显然成立。
\ans{必要性和充分性均得证；$I_1I_3<0$ 排除了点圆和虚圆。}

\Q{7}{直角双曲线的充要条件}
\textbf{题意：}两条渐近线相互垂直，证明等价于 $I_1=0,I_3\ne0$。\hint{以渐近线交点为中心，正交旋转到主轴形 $\lambda u^2+\mu v^2+k=0$。}
直角双曲线的两渐近方向满足 $\lambda u^2+\mu v^2=0$，其斜率分别为 $\pm\sqrt{-\lambda/\mu}$，互相垂直当且仅当 $\lambda=-\mu$，即 $I_1=\lambda+\mu=0$。非退化需 $I_3\ne0$。
反之 $I_1=0,I_3\ne0$ 意味着 $A$ 的非零特征值互为相反数（$I_2=-\lambda^2<0$），所以必是真双曲线；其渐近线与两主轴各成 $45^\circ$，两条渐近线正交。
\ans{$\boxed{I_1=0,\ I_3\ne0}$ 为所求充要条件。}
